% =====================================================================
% 5.3 실무적 기여 — 3차심사 대규모 재설계 (발견→도구 축)
%   5.3.1 회귀식 기반 DRTO 산출·예측 도구          [신설]
%   5.3.2 이중채널 감쇠 기반 자원배분·스트레스테스트 [신설]
%   5.3.3 η 기반 평상시 성숙도·위기시 에스컬레이션   [구 5.3.1 성숙도 + 구 5.3.3 에스컬레이션 통합]
%   5.3.4 산업별 적용 가이드라인                    [구 5.3.2]
% =====================================================================

\section{\vrev{실무적 기여}}
\label{sec:practical-contrib}
\label{sec:practical-implications}    % 구 라벨 호환 — 외부 cross-ref 유지

\begin{p5add}
\chg{본 절은 본 연구가 BCM 실무 현장에 새롭게 제공한 의사결정 도구를 네 항으로 정리한다. 핵심은 본 연구의
두 핵심 발견인 회귀식(간명 모형)과 이중채널 감쇠 메커니즘을 각각 독립된 실무 도구로 구체화하여, 이론적
발견이 현장의 산출과 자원 배분으로 직접 내려오는 경로를 제시하는 데 있다.}
\chg{먼저 회귀식 기반 $\DRTO$ 산출과 예측 도구(5.3.1항)에서는 간명 회귀 모형으로 조직이 자기 환경의 $\DRTO$를
직접 계산하는 절차를, 이중채널 감쇠 기반 자원배분과 스트레스 테스트(5.3.2항)에서는 P85.8 분할과 $0.52$배
감쇠를 자원 배분 규칙과 극단 시나리오 설계로 전환하는 방안을 제시한다. 이어 $\eta$ 기반 평상시 성숙도와
위기시 에스컬레이션(5.3.3항)에서는 단일 지표 $\eta$의 평상시 진단과 위기시 대응의 이중 활용 및 선순환
구조를, 산업별 적용 가이드라인(5.3.4항)에서는 6개 대표 산업의 권장 파라미터와 분포 환경별 적용 원칙을
제시한다. 이들은 모두 \chref{ch:results}의 시뮬레이션과 회귀분석, 전문가 검증에서 도출된 정량적 결과를
실무 운용 형태로 변환한 산출물이다.}
\end{p5add}


% ===================================================================
% 5.3.1 회귀식 기반 DRTO 산출·예측 도구 [신설]
% ===================================================================
\subsection{\chg{회귀식 기반 $\DRTO$ 산출과 예측 도구}}
\label{subsec:practical-regression-tool}

\chg{본 항은 \chref{ch:results}의 회귀분석에서 도출된 간명 회귀 모형을 실무 현장의 $\DRTO$ 산출과 예측 도구로
제시한다. 이는 본 연구의 핵심 산출물인 회귀식이 단순한 통계적 적합을 넘어, 조직이 자기 환경의 $\DRTO$를
직접 계산하는 운영 도구로 활용되는 경로를 구체화한 것이다.}

\chg{도구의 토대는 전진 선택법과 LASSO의 교차검증 합의로 도출된 4변수 간명 모형(식~(\vref{eq:c2_parsimonious}),
(\vref{eq:c3_parsimonious}))이다. 9변수 완전 2차 모형이 C2에서 $R^2 = 0.998$, C3에서 $R^2 = 0.858$의
설명력을 보이는 데 비해, 합의 4변수 간명 모형은 C2 $R^2 = 0.996$, C3 $R^2 = 0.839$로 설명력 손실을
최소화하면서 해석 가능성을 확보한다. 특히 C1과 C2의 2차교호작용 회귀가 각각 $R^2 = 1.000$, $0.998$에
도달한다는 결과는 개별 시그모이드 바운딩 함수가 다항식으로 거의 완전히 근사됨을 의미하며, 이는 회귀식
산출 결과의 신뢰성을 뒷받침한다.}

\chg{실무 산출 절차는 다음과 같다. 조직은 자기 환경의 네 입력값, 곧 기준선 복구시간 $R_0$, 관측성 원시값
$O$, 관측성 가중치 $k_O$, 불확실성 수준 $U$를 간명 회귀식에 대입하여 효율비 추정값 $\hat\eta$를 얻고,
$\DRTO = \hat\eta \cdot R_0$로 동적 복구목표시간을 즉시 산출한다. 이렇게 산출된 $\DRTO$를 기존 BIA에서
고정값으로 설정한 정적 RTO와 비교하면, 관측성 투자와 불확실성 환경이 복구 목표에 미치는 영향을 정량적으로
확인할 수 있다. \tabref{tab:ch6-drto-examples}\refpo{tab:ch6-drto-examples}{은}{는} 뒤에 나올 산업별 권장
파라미터(\tabref{tab:ch6-industry-params})를 회귀식에 대입한 구체적 산출 예시이다.}

\begin{table}[htbp]
\centering
\caption{\chg{산업별 대표 $\DRTO$ 산출 예시}}
\label{tab:ch6-drto-examples}
\small
\begin{tabular}{l c c c c c}
\toprule
\textbf{산업} & \textbf{$R_0$ (h)} & \textbf{$O$} & \textbf{$\hat{\eta}$} & \textbf{$\DRTO$ (h)} & \textbf{조건} \\
\midrule
제조업   & 4.0 & 0.60 & $1.664$ & $6.66$ & C2 \\
금융     & 1.0 & 0.78 & $1.569$ & $1.57$ & C3 \\
의료     & 2.0 & 0.68 & $1.569$ & $3.14$ & C3 \\
IT 서비스 & 0.5 & 0.83 & $1.599$ & $0.80$ & C2 \\
물류     & 6.0 & 0.53 & $1.684$ & $10.10$ & C2 \\
소매     & 3.0 & 0.58 & $1.669$ & $5.01$ & C2 \\
\bottomrule
\end{tabular}
\end{table}

\chg{산출에는 4변수 간명 모형(C2 $R^2 = 0.996$, C3 $R^2 = 0.839$)을 적용하였으며, $U$는 산업별 불확실성
대표값(중앙값), $O_{\mathrm{raw}} = 0.5$를 가정하였다. 예컨대 금융 산업의 $\DRTO = 1.57$h는 $R_0 = 1.0$h
대비 $57\%$의 안전 마진을 포함한 것으로 파레토 불확실성에 의한 꼬리 위험을 반영한 값이며, 물류 산업의
$\DRTO = 10.10$h는 $R_0 = 6.0$h 대비 $68\%$ 증가한 것으로 비교적 낮은 관측성($O = 0.53$)과 불확실성
프리미엄이 결합된 결과이다. 이처럼 실무자는 자기 조직의 파라미터를 대입하여 $\DRTO$를 즉시 산출하고
기존 정적 RTO와 비교함으로써 동적 복구목표시간의 실무적 유용성을 직관적으로 확인할 수 있다.}

\chg{나아가 회귀식의 변수 진입 순서는 산출을 넘어 투자 우선순위 신호로도 활용된다. 4장의 전진 선택법
결과, 가우시안 환경(C2)에서는 관측성 가중치 $k_O$가 1순위로 진입하여 관측성이 $\eta$ 변동의 주된 결정
요인인 반면, 파레토 환경(C3)에서는 불확실성 $U$가 1순위로 진입하여 지배 변수가 전환된다. 따라서 조직은
자기 환경의 분포 특성에 따라, 가우시안 환경에서는 관측성 투자를, 파레토 환경에서는 불확실성 관리를
우선해야 한다는 자원 배분 신호를 회귀식으로부터 직접 읽어낼 수 있다.}

\chg{다만 간명 회귀식은 평균 영역의 거동을 근사하는 모형이므로, 극단 불확실성이 지배하는 꼬리 영역에서는
산출값의 정밀도가 제한된다. 꼬리 영역에서의 비선형 감쇠와 그에 따른 자원 배분은 다음 항(\subsecref{subsec:practical-dualchannel-tool})의
이중채널 도구에서 다룬다.}


% ===================================================================
% 5.3.2 이중채널 감쇠 기반 자원배분·스트레스테스트 [신설]
% ===================================================================
\subsection{\chg{이중채널 감쇠 기반 자원배분과 스트레스 테스트}}
\label{subsec:practical-dualchannel-tool}

\chg{본 항은 4장에서 발견한 이중채널 감쇠 메커니즘을 실무의 자원 배분 의사결정과 스트레스 테스트 설계
도구로 전환한다. 이중채널 감쇠란 P85.8 분할점을 경계로 관측성 가중치 $k_O$의 회귀 계수가 핵심 영역(Core,
$-0.798$)에서 꼬리 영역(Tail, $-0.415$)으로 $0.52$배 감쇠하는 현상으로, 극단 불확실성 환경에서 시그모이드
포화에 의해 관측성 투자의 한계 효용이 줄어든다는 것을 정량화한 본 연구의 핵심 발견이다.}

\chg{첫째, P85.8 분할점은 조직이 자기 불확실성 환경이 어느 체제에 있는지를 진단하고 대응 모드를 전환하는
기준이 된다. 가우시안(C2)과 파레토(C3) 분포의 누적분포함수가 P85.8($\eta \approx 2.42$)에서 교차하므로,
이 지점 이하의 일상 운영(Core)에서는 두 환경이 사실상 구별되지 않으나(Cohen's $d \approx 0$), 이 지점을
넘어서면 파레토 꼬리가 급격히 발현하여 $\mathrm{CVaR}_{99}$ 기준 C3가 C2를 $+24.0\%$ 상회한다. 따라서
조직의 복구 시간 분포가 P85.8 분할점을 넘어서기 시작하는 시점, 곧 가우시안 환경에서 파레토 환경으로의
전환은 그 자체로 자원 배분을 재조정해야 할 신호로 취급된다.}

\chg{둘째, 이중채널 감쇠 비율은 영역별 자원 배분 규칙의 정량적 근거를 제공한다. Core 영역에서는 관측성
가중치를 $0.1$ 높이면 $\eta$가 $0.080$ 감소하여 관측성 투자의 한계 수익이 충분히 발현되므로, 일상 운영의
복구 역량 향상은 관측성 투자(모니터링 확대, APM/SIEM 연동)에 집중하는 것이 효율적이다. 반면 Tail 영역에서는
동일한 관측성 개선이 $\eta$를 $0.042$만 감소시켜 한계 수익이 Core의 $0.52$배로 감쇠하므로, 극단 상황에서는
관측성 투자만으로 복구 시간을 충분히 단축하기 어렵다. 이 영역에서는 불확실성 자체의 관리, 곧
업무연속성계획(BCP) 시나리오의 다양화, 대체 사이트 확보, 사이버 보험과 같은 꼬리 위험 전가 수단에 자원을
전환해야 한다. 이는 동일한 예산이라도 환경 체제에 따라 투자 대상이 질적으로 달라져야 함을 의미한다.}

\chg{셋째, 이 발견은 스트레스 테스트 설계의 기준을 제시한다. $0.52$배 감쇠는 평균 영역의 시나리오만으로는
극단 환경에서의 복구 역량을 검증할 수 없음을 뜻하므로, BCM 스트레스 테스트는 Tail 영역의 시나리오를 별도로
설계해야 한다. 구체적으로, 파레토 분포의 P99 수준 극단값을 입력으로 하는 시나리오에서 관측성 투자만으로는
목표 복구 시간을 달성하지 못함을 검증하고, 불확실성 관리 수단을 병행했을 때의 효과를 함께 측정하는 이중
시나리오 설계가 요구된다.}

\chg{넷째, 환경 전환에 따른 회귀 계수 변화는 모델 재교정의 신호로 활용된다. 4장의 회귀계수 직접 비교 결과,
가우시안에서 파레토로 환경이 전환되면 관측성 변수 $k_O$의 계수가 $0.75$배, 불확실성 변수 $U$의 선형 계수가
$0.45$배로 변화한다. 또한 관측성 투자의 순효과는 가우시안 경로(Cohen's $d = +1.110$)보다 파레토 경로(Cohen's
$d = +0.650$)에서 상대적으로 작게 나타나, 환경에 따라 관측성 투자의 순효과 크기가 달라진다는 점도 확인된다.
조직은 자기 환경의 분포가 전환되는 것을 감지하면 회귀식의 계수를 재추정하여 산출 도구(\subsecref{subsec:practical-regression-tool})를
갱신해야 한다.}

\chg{요컨대 이중채널 감쇠 메커니즘은 ``일상에서는 관측성 투자, 극단에서는 불확실성 관리''라는 이중축 자원
배분 원칙과, Tail 시나리오를 별도로 설계하는 스트레스 테스트 기준을 동시에 제공한다. 이는 단일 지표나 평균
중심 의사결정으로는 포착할 수 없는, 분포 환경에 따른 입체적 위기 대응을 가능하게 하는 본 논문의 독자적
실무 기여이다.}


% ===================================================================
% 5.3.3 η 기반 평상시 성숙도·위기시 에스컬레이션 [구 성숙도 + 에스컬레이션 통합]
% ===================================================================
\subsection{\chg{$\eta$ 기반 평상시 성숙도와 위기시 에스컬레이션}}
\label{subsec:practical-bcm-std}        % 구 라벨 호환(성숙도)
\label{subsec:practical-decision}       % 구 라벨 호환(의사결정)
\label{subsec:escalation-thresholds}    % 구 라벨 호환
\label{subsec:eta-dual-use}             % 구 라벨 호환

\begin{p5add}
\chg{본 항은 단일 지표 $\eta = \DRTO/R_0$가 운영 모드에 따라 평상시 성숙도 진단 지표와 위기시 에스컬레이션
대응 지표의 두 역할을 수행함을 보이고, 그 평상시와 위기시 활용 및 선순환 구조를 제시한다. 평상시에는 $\eta$
변동이 주로 관측성 채널에 의해 결정되므로 BCMS 성숙도를 5단계 모형으로 진단하고, 위기시에는 $\eta$ 변동이
불확실성 채널에 의해 지배되므로 Green/Yellow/Orange/Red 4단계 에스컬레이션을 자동 트리거한다. 또한 데이터
수집과 모니터링 대시보드, 사후 분석으로 구성되는 BCM 운영 통합 도구가 함께 제공된다.}
\end{p5add}

\begingroup
  \renewcommand{\section}[1]{}
  \renewcommand{\subsection}[1]{}
  \renewcommand{\paragraph}[1]{\textbf{#1}\par}

  % η의 평상시·위기시 이중 활용 개념(overview) + 성숙도 5단계(1.1) + BCM 운영통합(1.2-4)
  
$\eta = \DRTO / R_0$는 단일 지표이지만, 조직의 운영 상황에 따라
두 가지 서로 다른 역할을 수행한다.
이러한 이중 활용이 가능한 근거는 $\eta$를 구성하는 두 채널\chg{(}관측성($E_{O,bnd}$)과
불확실성($E_{U,bnd}$)\chg{)}의 \textbf{지배적 기여 요인이 운영 모드에 따라 달라지기}
때문이다.
평상시에는 \chg{대규모 재난과 사고}가 부재하므로 불확실성이 \chg{일상적이고 가우시안}
수준에 머물러 $\eta$의 변동을 주로 관측성 역량이 결정한다.
반면 위기시에는 이미 \chg{사고와 재난}이 발생하여 규모를 알 수 없는 큰 불확실성이
현존하며, 이를 얼마나 빠르게 정량화하여 $\eta$에 반영하느냐가 대응의 핵심이 된다.
\p{[}표~\vref{tab:eta-dual-use}\p{]}는 이 \chg{평상시와 위기시} 활용 체계를 정리한 것이다.

\begin{table}[htbp]
\centering
\caption{$\eta$ 지표의 \chg{평상시·위기시} 활용 체계}
\label{tab:eta-dual-use}
\small
\begin{tabularx}{\textwidth}{l X X}
\toprule
\textbf{구분} & \textbf{평상시 (Normal Mode)} & \textbf{위기시 (Crisis Mode)} \\
\midrule
\textbf{역할} & BCMS 성숙도 평가 지표 & 에스컬레이션 대응 지표 \\
\textbf{핵심 질문} & ``조직의 복구 역량이 어느 수준인가?'' & ``현재 상황이 얼마나 심각한가?'' \\
\textbf{시간축} & 중장기 (월간/분기별 추이) & 실시간 (분$\sim$시간 단위) \\
\textbf{주요 활동} & 진단, 투자 계획, 벤치마크 & 등급 판정, 보고, 즉각 대응 \\
\textbf{대상자} & BCM 관리자, 경영진 & 복구 팀, 현장 담당자 \\
\textbf{산출물} & 성숙도 등급, 개선 로드맵 & 에스컬레이션 등급, 대응 조치 \\
\bottomrule
\end{tabularx}
\end{table}

\p{[}그림~\vref{fig:eta-dual-use-diagram}\p{]}는 $\eta$의 \chg{평상시와 위기시} 이중 활용과
BCMS 선순환 구조를 종합적으로 도식화한 것이다.
평상시 성숙도 평가(좌측)와 위기시 에스컬레이션 대응(우측)이
\chg{장애 발생과 복구 완료}를 통해 전환되며,
사후 분석과 역량 환류가 상단의 진단--향상--대응--환류 선순환을 형성한다.

\begin{figure}[htbp]
\centering
\resizebox{\textwidth}{!}{%
\begin{tikzpicture}[
  arr/.style={-{Stealth[length=2.5mm, width=2mm]}, line width=1.2pt},
  arrlbl/.style={font=\scriptsize\bfseries, fill=white, inner sep=1.5pt, rounded corners=1.5pt},
  sbox/.style={draw, rounded corners=2pt, semithick, minimum width=20mm,
    minimum height=9mm, align=center, font=\scriptsize},
]

%% ============================================================
%% 상단 박스: BCMS 선순환
%% ============================================================
\node[draw, rounded corners=6pt, line width=1pt,
  fill=green!6, draw=green!50!black,
  minimum width=120mm, minimum height=24mm] (topbox) at (0, 3.2) {};

\node[font=\small\bfseries, text=black] at (0, 4.15) {BCMS 선순환};

\node[sbox, fill=blue!12, draw=blue!50]  (s1) at (-3.7, 3.2) {\textbf{1. 진단}\\$\eta$ 추이};
\node[sbox, fill=blue!12, draw=blue!50]  (s2) at (-1.25, 3.2) {\textbf{2. 향상}\\$O\!\uparrow\; U\!\downarrow$};
\node[sbox, fill=red!12, draw=red!50]    (s3) at (1.25, 3.2)  {\textbf{3. 대응}\\에스컬레이션};
\node[sbox, fill=yellow!18, draw=yellow!60!black] (s4) at (3.7, 3.2) {\textbf{4. 환류}\\이력$\to$KPI};

\draw[-{Stealth[length=1.8mm]}, thin, blue!60] (s1) -- (s2);
\draw[-{Stealth[length=1.8mm]}, thin, red!50] (s2) -- (s3);
\draw[-{Stealth[length=1.8mm]}, thin, yellow!60!black] (s3) -- (s4);
\draw[-{Stealth[length=1.8mm]}, thin, green!50!black]
  (s4.south) -- ++(0,-0.28) -| (s1.south)
  node[pos=0.5, below, font=\scriptsize\bfseries, text=black,
    fill=green!6, inner sep=1.5pt, rounded corners=1.5pt] {$\eta$ 선순환};

%% ============================================================
%% 좌하 박스: 평상시 BCMS 성숙도
%% ============================================================
\node[draw, rounded corners=6pt, line width=1pt,
  fill=blue!6, draw=blue!60,
  minimum width=56mm, minimum height=38mm] (leftbox) at (-5.0, -1.0) {};

\node[font=\small\bfseries, text=black] at (-5.0, 0.55) {평상시: BCMS 성숙도 평가};

\node[anchor=center, inner sep=0pt] at (-5.0, -0.85) {%
  \scalebox{0.7}{%
  \begin{tabular}{c c c l}
  \toprule
  \textbf{단계} & \textbf{등급} & \textbf{$\eta$} & \textbf{수준} \\
  \midrule
  5 & 최적화 & $< 0.6$ & \textit{Optimized} \\
  4 & 관리 & $0.6$--$0.8$ & \textit{Managed} \\
  3 & 정의 & $0.8$--$1.0$ & \textit{Defined} \\
  2 & 반복 & $1.0$--$1.5$ & \textit{Repeatable} \\
  1 & 초기 & $\geq 1.5$ & \textit{Initial} \\
  \bottomrule
  \end{tabular}}%
};

\node[font=\scriptsize\bfseries, text=black] at (-5.0, -2.35) {%
  개선 두 축: 관측성$\uparrow$ / 불확실성$\downarrow$%
};

%% ============================================================
%% 우하 박스: 위기시 에스컬레이션
%% ============================================================
\node[draw, rounded corners=6pt, line width=1pt,
  fill=red!6, draw=red!60,
  minimum width=56mm, minimum height=38mm] (rightbox) at (5.0, -1.0) {};

\node[font=\small\bfseries, text=black] at (5.0, 0.55) {위기시: 에스컬레이션 대응};

\node[anchor=center, inner sep=0pt] at (5.0, -0.85) {%
  \scalebox{0.7}{%
  \begin{tabular}{c c l l}
  \toprule
  \textbf{$\eta$} & \textbf{등급} & \textbf{보고} & \textbf{대응} \\
  \midrule
  $< 0.8$ & Green & 팀 자체 & 정상 복구 \\
  $0.8$--$1.0$ & Yellow & 팀장 & 모니터링 강화 \\
  $1.0$--$1.5$ & Orange & 부서장 & 자원 추가 투입 \\
  $\geq 1.5$ & Red & 경영진 & BCP 전면 가동 \\
  \bottomrule
  \end{tabular}}%
};

\node[font=\scriptsize\bfseries, text=black] at (5.0, -2.35) {%
  국가 위기경보 연동: 관심--주의--경계--심각%
};

%% ============================================================
%% 3개 박스 연결 화살표
%% ============================================================
\draw[arr, red!50]
  ([yshift=-4mm]leftbox.east) -- ([yshift=-4mm]rightbox.west)
  node[arrlbl, midway, below=1.5pt, text=black] {장애$\cdot$재난 발생 $\to$ 위기 모드 전환};
\draw[arr, blue!50]
  ([yshift=4mm]rightbox.west) -- ([yshift=4mm]leftbox.east)
  node[arrlbl, midway, above=1.5pt, text=black] {복구 완료 $\to$ 평상 모드 복귀};
\draw[arr, green!50!black]
  ([xshift=6mm]rightbox.north) -- ([xshift=17mm]topbox.south)
  node[arrlbl, midway, right=1.5pt, text=black] {사후 분석};
\draw[arr, blue!60]
  ([xshift=-17mm]topbox.south) -- ([xshift=-6mm]leftbox.north)
  node[arrlbl, midway, left=1.5pt, text=black] {역량 환류};

%% ============================================================
%% 하단 강조 문구
%% ============================================================
\node[font=\small\bfseries, text=black, align=center,
  draw=red!50, rounded corners=3pt, fill=red!6, inner sep=4pt,
  line width=0.6pt] at (0, -3.5) {%
  $\eta$ 하나로 평상시 ``얼마나 준비되었는가'' + 위기시 ``얼마나 심각한가'' 통합 측정
  $\;\to\;$ \textcolor{black}{\bfseries 데이터 기반 BCMS 선순환}%
};

\end{tikzpicture}%
}% end resizebox
\caption{$\eta$의 \chg{평상시와 위기시} 이중 활용 체계와 BCMS 선순환 구조}
\label{fig:eta-dual-use-diagram}
\end{figure}

  %% --- 평상시: BCMS 성숙도 평가 ---
\subsubsection{\tmp{평상시 BCMS 성숙도 평가}}
\label{subsubsec:maturity-model}

평상시에는 \chg{대규모 사고와 재난}이 발생하지 않은 상태이므로,
불확실성은 일상적 업무 변동(가우시안 수준)에 한정되어 $E_{U,bnd}$의 기여가 작다.
따라서 $\eta$의 변동은 주로 \textbf{관측성 역량}($E_{O,bnd}$)에 의해 결정되며,
$\eta$의 중장기 추이를 추적하면 조직이 모니터링 체계를 얼마나 잘 구축하고
있는지\chg{,} 즉 복구 역량의 성숙도를 정량적으로 진단할 수 있다.
$\eta$가 낮을수록 관측성 역량이 우수한 성숙한 조직으로 해석되며,
이에 기반하여 BCMS 성숙도를 5단계로 구분하는 모형을 \p{[}표~\vref{tab:maturity-5level}\p{]}에서 제안한다.

각 임계값은 본 연구의 시뮬레이션 결과에 기반한다\chg{.}
$\eta = 0.6$은 C1 조건의 하위 \erf{25}{5}\% 수준으로 관측성이 극대화된 상태,
$\eta = 0.8$은 C1 평균 $\bar{\eta} = 0.853$ 이하로 $\DRTO < R_0$가 안정적으로 유지되는 상태,
$\eta = 1.0$은 $\DRTO = R_0$ 기준선으로 관측성과 불확실성이 상쇄되는 지점,
$\eta = 1.5$는 불확실성이 지배적이어서 $\DRTO$가 $R_0$의 1.5배를 초과하는 상태이다.

\begin{table}[htbp]
\centering
\caption{$\eta$ 기반 BCMS 성숙도 5단계 모형}
\label{tab:maturity-5level}
\small
\begin{tabularx}{\textwidth}{c c c l X}
\toprule
\textbf{단계} & \textbf{등급} & \textbf{$\eta$ 범위} & \textbf{수준} & \textbf{특징} \\
\midrule
5 & 최적화 & $\eta < 0.6$ & \textit{Optimized} &
  관측성 극대화, 불확실성 최소화. 자동화된 모니터링·복구 \\
\addlinespace
4 & 관리 & $0.6 \leq \eta < 0.8$ & \textit{Managed} &
  정량적 KPI 기반 BCM 운영. $\DRTO < R_0$ 안정 유지 \\
\addlinespace
3 & 정의 & $0.8 \leq \eta < 1.0$ & \textit{Defined} &
  공식 BCM 프로세스 수립. 관측성과 불확실성 \m{상쇄($\DRTO \approx R_0$)} \\
\addlinespace
2 & 반복 & $1.0 \leq \eta < 1.5$ & \textit{Repeatable} &
  기본 절차 존재하나 불확실성 우세. $\DRTO > R_0$ 빈번 \\
\addlinespace
1 & 초기 & $\eta \geq 1.5$ & \textit{Initial} &
  BCM 체계 미비. 극단적 불확실성 노출 \\
\bottomrule
\end{tabularx}
\end{table}

성숙도를 향상(즉, $\eta$를 감소)시키는 경로는 두 축으로 구분된다.
\chg{첫째, 관측성 강화($E_{O,bnd}$ 증가로 $\eta$ 감소)는 실시간 모니터링 체계 구축, APM/SIEM 연동, 장애
원인 분석(RCA) 체계화를 포함하며, 본 연구의 C1 조건에서 $k_O$ 계수가 관측성의 실질적 효과를 결정하므로
$k_O$가 높은 영역에 우선 투자하는 것이 효율적이다. 둘째, 불확실성 저감($E_{U,bnd}$ 감소로 $\eta$ 감소)은
복구 절차서 최신 유지, 정기 BCM 훈련, 공급망 대체 방안 확보를 포함하며, 특히 Tail 영역(P85.8 초과)에서는
파레토 불확실성이 지배적이므로 블랙스완 대비 계획이 성숙도 향상의 핵심이 된다.}

각 성숙도 단계별 구체적 전환 전략, 위기 시 정책 활용 방안, 도입 로드맵은
아래에 상세히 기술하며, 현업 담당자용 자가 진단 체크리스트(15개 항목)는
\appchref{app:eta-dual-use-detail}에 수록한다.

%% [분량보존 복원] η 평상시·비상시 활용 상세 (구 부록 H.2~H.4): 향상전략·위기정책·로드맵
\subsubsection{성숙도 단계별 향상 전략}
\label{app:maturity-strategy}

각 성숙도 단계에서 다음 단계로 향상하기 위한 구체적 전략을 기술한다.
핵심은 관측성 강화($E_{O,bnd}$ 증가)와 불확실성 저감($E_{U,bnd}$ 감소)의
두 레버를 단계에 맞게 적용하는 것이다.

\chg{1단계(Initial)에서 2단계(Repeatable)로의 기반 구축 단계에서는 기본 모니터링 도구 도입(가용성
모니터링, 기본 알림), BCM 정책 문서 및 복구 절차서(runbook) 초안 작성, BIA 최초 수행을 통한 RTO
기준값($R_0$) 설정, RPO 기준값 설정, BCM 전담 책임자 지정을 수행하며, 목표는 $\eta < 1.5$ 안정화이다.}

\chg{2단계(Repeatable)에서 3단계(Defined)로의 체계화 단계에서는 관측성 3축(로그, 메트릭, 트레이스) 중
2개 이상 확보, 장애 대응 절차 표준화 및 정기 훈련(연 2회 이상), MTTR 측정과 추적 시작을 통한
$O_{\text{raw}}$ 데이터 수집, 에스컬레이션 기준 및 보고 체계 문서화를 수행하며, 목표는 $\eta < 1.0$($\DRTO \leq R_0$)이다.}

\chg{3단계(Defined)에서 4단계(Managed)로의 정량화 단계에서는 실시간 $\DRTO/R_0$ 모니터링을 위한 $\eta$
대시보드 구축, $k_O$ 변화 추적을 통한 관측성 투자 ROI 측정, 불확실성 원인별 분류 체계 수립(가우시안 vs
파레토 유형), 사후 분석을 통한 절차 개선 피드백 루프 정착을 수행하며, 목표는 $\eta < 0.8$ 유지(에스컬레이션
Green 영역)이다.}

\chg{4단계(Managed)에서 5단계(Optimized)로의 최적화 단계에서는 AIOps 기반 자동 장애 탐지와 진단, 복구
체계 구축, 카오스 엔지니어링(Chaos Engineering) 정기 수행, 향후 $\eta$ 추이 예측 및 선제 대응을 위한
$\eta$ 예측 모형 운영, 극단 시나리오(C3) 전용 대응을 위한 CVaR$_{99}$ 기반 자원 배분을 수행하며, 목표는
$\eta < 0.6$ 달성 및 유지이다.}


\subsubsection{위기시 정책 활용 방안}
\label{app:crisis-policy}

\chg{국가 위기경보 연동 $\eta$ 임계값 자동 조정은}
국가 위기경보가 상향될 때 조직 내부의 에스컬레이션 민감도를
자동으로 높여 선제 대응하는 정책이다.
\nref{위기경보 4단계에 따른 $\eta$ 임계값 조정 규칙은 \tabref{tab:app-policy-adjustment}와 같다.}

\begin{table}[htbp]
\centering
\caption{국가 위기경보 연동 $\eta$ 임계값 조정 정책}
\label{tab:app-policy-adjustment}
\small
\newcolumntype{W}{>{\raggedright\arraybackslash}p{0.42\textwidth}}
\begin{tabularx}{\textwidth}{c c X W}
\toprule
\textbf{위기경보} & \textbf{임계값 조정} & \textbf{근거} & \textbf{조정 후 기준} \\
\midrule
관심(Blue) & 기본값 유지 &
  일상 운영 &
  Green $\eta < 0.8$,\; Yellow $0.8 \leq \eta < 1.0$,\; Orange $1.0 \leq \eta < 1.5$,\; Red $\eta \geq 1.5$ \\
\addlinespace
주의(Yellow) & 각 $-0.1$ &
  외부 위험 증가 &
  Green $\eta < 0.7$,\; Yellow $0.7 \leq \eta < 0.9$,\; Orange $0.9 \leq \eta < 1.4$,\; Red $\eta \geq 1.4$ \\
\addlinespace
경계(Orange) & 각 $-0.2$ &
  위기 가능성 높음 &
  Green $\eta < 0.6$,\; Yellow $0.6 \leq \eta < 0.8$,\; Orange $0.8 \leq \eta < 1.3$,\; Red $\eta \geq 1.3$ \\
\addlinespace
심각(Red) & 즉시 Red &
  대규모 위기 &
  $\eta$와 무관하게 BCP 전면 가동, 전 등급 Red 적용 \\
\bottomrule
\end{tabularx}
\end{table}


\chg{위기 유형별 $\eta$ 가중 모형은}
재난 유형에 따라 $\eta$ 구성 요소(관측성 vs 불확실성)에 가중치를 부여한다.
\nref{재난 4유형별 $O\!:\!U$ 가중 비율과 대응 전략 초점은 \tabref{tab:app-crisis-type-weight}에 정리하였다.}

\begin{table}[htbp]
\centering
\caption{위기 유형별 대응 전략 가중 모델}
\label{tab:app-crisis-type-weight}
\small
\begin{tabularx}{\textwidth}{l c c X}
\toprule
\textbf{위기 유형} & \textbf{$O$ 비중} & \textbf{$U$ 비중} & \textbf{대응 전략 초점} \\
\midrule
사이버 공격 & 높음 (70\%) & 보통 (30\%) &
  SIEM/EDR 강화, 탐지·진단 역량 집중 \\
\addlinespace
자연재해 & 보통 (40\%) & 높음 (60\%) &
  물리적 피해 불확실성 대비, 대체 사이트 확보 \\
\addlinespace
감염병 & 보통 (50\%) & 높음 (50\%) &
  원격 운영 역량 + 인력 가용성 불확실성 동시 대응 \\
\addlinespace
공급망 중단 & 낮음 (30\%) & 높음 (70\%) &
  외부 불확실성 지배. 다중 공급원 확보, 재고 완충 \\
\bottomrule
\end{tabularx}
\end{table}


\chg{에스컬레이션 이력 기반 성숙도 KPI로서,}
위기시 에스컬레이션 이력은 평상시 성숙도 개선의 데이터로 환류된다.
\chg{주요 KPI는 네 가지이다. Green 체류 비율은 전체 운영 시간 대비 $\eta < 0.8$ 비율로 성숙도가 높을수록
증가하며, 에스컬레이션 빈도는 월간 Yellow 이상 발생 횟수로 감소 추세가 성숙도 향상을 의미하고, 복귀
시간은 Orange/Red에서 Green으로의 복귀 소요 평균 시간으로 단축이 복구 역량 향상을 뜻하며, Red 도달률은
위기(Red) 등급 발생 빈도로 감소가 극단 대비 역량 향상을 나타낸다.}

이러한 KPI를 월간/분기별로 경영진에 보고함으로써
BCM 투자의 효과를 정량적으로 증명할 수 있다.


\subsubsection{$\eta$ \chg{평상시와 비상시} 활용 체계 도입 로드맵}
\label{app:eta-roadmap}

\nref{$\eta$ 기반 BCM 체계의 단계별 도입 일정과 주요 활동은 \tabref{tab:app-roadmap}에 정리하였다.}

\begin{table}[htbp]
\centering
\caption{$\eta$ 평상시·비상시 활용 체계 도입 로드맵}
\label{tab:app-roadmap}
\small
\begin{tabularx}{\textwidth}{c l X}
\toprule
\textbf{단계} & \textbf{기간} & \textbf{주요 활동} \\
\midrule
준비 & 1--2개월 &
  BIA 수행 $\to$ $R_0$ 설정. 관측성 데이터 수집 시작.
  $\eta$ 산출 파일럿 구축. 자가 진단 최초 실시 \\
\addlinespace
파일럿 & 3--4개월 &
  핵심 시스템 1--2개에 $\eta$ 대시보드 적용.
  에스컬레이션 기준 시범 운영. 국가 위기경보 연동 테스트 \\
\addlinespace
확산 & 5--8개월 &
  전사 확대 적용. 성숙도 자가 진단 분기별 정기화.
  에스컬레이션 이력 KPI 보고 시작 \\
\addlinespace
최적화 & 9--12개월 &
  AIOps 연동. $\eta$ 예측 모형 운영.
  위기 유형별 가중 모델 적용. 산업별 벤치마크 비교 \\
\bottomrule
\end{tabularx}
\end{table}






  %% -------------------------------------------------------------------
\subsection{\vrev{BCM 운영 통합 프레임워크}}
\label{subsec:bcm-integration}
%% -------------------------------------------------------------------

2단계 프레임워크와 에스컬레이션 체계를 기존 BCMS에 통합하기 위한
구체적 운영 절차를 다음과 같이 제시한다.

\subsubsection{\tmp{데이터 수집 체계}}

\inew{$\DRTO$ 모델의 두 입력 변수인 \textbf{관측성 $O(t)$} 와 \textbf{불확실성 $U(t)$} 의
수집은 각각 기존 \chg{운영과 보안} 모니터링 인프라와 \chg{장애 이력 및 위협 인텔리전스} 데이터의
연동을 통해 구현한다.}

\inew{\chg{관측성 지표 $O(t)$의 수집은 기존 모니터링 인프라와의 연동을 통해 다음과 같이 구현한다. 첫째, APM(Application Performance Monitoring) 연동으로 서비스 응답 시간, 오류율, 처리량을 종합하여 $O_{\mathrm{APM}} \in [0,1]$로 정규화하며, 가중 평균은 $O_{\mathrm{APM}} = w_1 \cdot R_{\mathrm{avail}} + w_2 \cdot (1-R_{\mathrm{error}}) + w_3 \cdot R_{\mathrm{perf}}$($\sum w_i = 1$)로 산출한다. 둘째, SIEM(Security Information and Event Management) 연동으로 보안 이벤트 심각도와 대응 상태를 기반으로 $O_{\mathrm{SIEM}} \in [0,1]$을 산출하되, 미대응 고위험 이벤트가 존재하면 감점한다. 셋째, 종합 관측성은 $O(t) = \alpha \cdot O_{\mathrm{APM}} + (1-\alpha) \cdot O_{\mathrm{SIEM}}$로 정의하며, $\alpha$는 조직 특성에 따라 예를 들어 $[0.5, 0.8]$에서 설정한다.}}

\inew{\chg{불확실성 지표 $U(t)$의 수집은 조직 내부의 과거 장애 이력과 외부 위협 인텔리전스 데이터의 연동을 통해 다음과 같이 구현한다. 첫째, 장애 이력 DB 연동으로 조직 자체의 과거 장애와 복구 이력(MTTR, RTO 초과 사례, 복구 소요 시간 분포)을 적합 분포(가우시안 또는 파레토)로 모델링하여 $U_{\mathrm{hist}} \in [0,\; U_{\max}]$을 산출하며, 분포 형태 판정에는 Hill 추정량~\citep{hill1975simple} 또는 Anderson-Darling 검정~\citep{anderson1954test}을 적용하여 C2(경꼬리) 또는 C3(중꼬리) 환경으로 자동 분기한다. 둘째, 위협 인텔리전스(CVE/MITRE ATT\&CK) 연동으로 자산 대장과 매핑된 CVE 공개 DB, EPSS 점수, MITRE ATT\&CK 위협 액터 활동 정보를 기반으로 미래 위험 강도 $U_{\mathrm{threat}} \in [0,\; U_{\max}]$을 산출하되, 미패치 고위험 CVE 수와 EPSS 가중 평균을 결합한다. 셋째, 종합 불확실성은 $U(t) = \beta \cdot U_{\mathrm{hist}} + (1-\beta) \cdot U_{\mathrm{threat}}$로 정의하며, $\beta$는 조직의 과거 데이터 풍부도와 위협 환경 변동성에 따라 $[0.5, 0.8]$에서 설정한다. 장애 이력이 풍부한 성숙 조직은 $\beta$가 높고, 신생 조직과 신규 위협 환경은 $\beta$가 낮다.}}

\inew{원시 변수 $O_{\mathrm{raw}}$, $U_{\mathrm{raw}}$, 가중치 $k_O$의 제1계층
하위 지표 조작화 예시와 가상 기업 라이프사이클 적용 시나리오는
\secref{app:operationalization}(관측성 $O_{\mathrm{raw}}$ 조작화 예시,
불확실성 $U_{\mathrm{raw}}$ 조작화 예시)에 참고용으로 수록되어 있다.}

\subsubsection{\tmp{$\eta$ 모니터링 대시보드}}

$\eta$ 비율의 실시간 시각화를 통해 BCMS 효과성을 지속적으로 측정한다.
대시보드의 핵심 구성요소는 \chg{실시간 $\eta$ 추이 차트(시계열),
에스컬레이션 등급 표시(Green/Yellow/Orange/Red),
$O(t)$와 $U(t)$의 개별 추이, CVaR$_{99}$ 현재값}이며,
이는 \nrev{ISO 22301:2019} 조항~9.1(모니터링, 측정, 분석, 평가)의
정량적 요구사항을 직접 충족한다.

\subsubsection{\tmp{사후 분석 프로세스}}

복구 완료 후 24--72시간 이내에 다음의 사후 분석을 수행한다.
\chg{먼저 실제 복구 시간 $T_{\mathrm{actual}}$과 $\DRTO$ 예측값을 비교하고,
절대 오차 $E = T_{\mathrm{actual}} - \DRTO_{\mathrm{final}}$과
상대 오차 $E_r = E / T_{\mathrm{actual}}$를 산출하며,
과대 및 과소 추정 원인을 분석하여 파라미터 교정을 권고하고,
장애 이력 DB 갱신 및 $U(t)$ 분포 재추정을 수행한다.}
이 프로세스는 \nrev{ISO 22301:2019} 조항~10.1(지속적 개선)과 직접 대응된다.




  % 위기시 에스컬레이션(3.1) + 선순환(3.2)
  %% --- 위기시: 에스컬레이션 대응 체계 ---
\subsubsection{\syncr{위기 시 에스컬레이션 대응 체계}}
\label{subsubsec:escalation-detail}

위기시에는 \chg{사고와 재난}이 이미 발생한 상태이므로,
규모를 알 수 없는 큰 불확실성이 현존한다.
이때 $\eta$는 이 불확실성의 크기를 실시간으로 반영하는 위험 수준 지표가 되며,
조직이 불확실성을 얼마나 빠르게 정량화하여 $\eta$에 반영하느냐가
적시 대응의 관건이 된다.
$\eta$ 값에 따라 4단계 에스컬레이션 체계가 자동 결정된다.
\p{[}표~\vref{tab:ch6-eta-escalation}\p{]}는 $\eta$ 구간별 대응 체계를 정의하며,
\p{[}그림~\vref{fig:ch6-escalation-zones}\p{]}는 이를 시각화한 것이다.

\begin{table}[htbp]
\centering
\caption{$\eta$ 기반 에스컬레이션 임계값 및 대응 체계}
\label{tab:ch6-eta-escalation}
\small
\begin{tabularx}{\textwidth}{c c l l X}
\toprule
\textbf{단계} & \textbf{$\eta$ 범위} & \textbf{등급} & \textbf{보고 대상} & \textbf{대응 조치} \\
\midrule
1 & $\eta < 0.8$ & \m{안정(Green)} & 복구 팀 자체 & 정상 복구 절차 수행; 관측성에 의한 $\DRTO$ 단축 효과 활성 \\
\addlinespace
2 & $0.8 \leq \eta < 1.0$ & \m{주의(Yellow)} & 팀장 보고 & 불확실성 증가 모니터링; $\DRTO$ 갱신 주기 \m{단축(15분)}; 추가 자원 대기 \\
\addlinespace
3 & $1.0 \leq \eta < 1.5$ & \m{경고(Orange)} & 부서장 보고 & $\DRTO > R_0$로 기준 RTO 초과; 추가 자원 투입; 외부 지원 요청; 이해관계자 긴급 통보 \\
\addlinespace
4 & $\eta \geq 1.5$ & \m{위기(Red)} & 경영진 보고 & 위기관리 모드 전환; 사업 연속성 비상 계획(BCP) 전면 가동; CVaR$_{99}$ 기준 최악 시나리오 대비 \\
\bottomrule
\end{tabularx}
\end{table}

\chg{에스컬레이션 트리거는 $\eta$ 급등(30\% 이상 증가), 조건 상향 전환(C2 $\to$ C3), 에스컬레이션 임계값
접근($\eta > 1.4$), 불확실성 급등($U > 0.8$)의 네 가지이다.}

\begin{figure}[htbp]
\centering
\resizebox{0.95\textwidth}{!}{%
\begin{tikzpicture}
% 배경 영역
\fill[green!15]  (0,0)   rectangle (3.2,3.5);
\fill[yellow!20] (3.2,0) rectangle (6.4,3.5);
\fill[orange!20] (6.4,0) rectangle (11.2,3.5);
\fill[red!20]    (11.2,0) rectangle (14.4,3.5);

% 경계선
\draw[thick, black] (3.2,0) -- (3.2,3.5);
\draw[thick, black] (6.4,0) -- (6.4,3.5);
\draw[thick, black] (11.2,0) -- (11.2,3.5);

% 축
\draw[thick, -{Stealth[length=3mm]}] (0,0) -- (15,0) node[right, font=\small] {$\eta$};
\draw[thick] (0,0) -- (0,3.5);

% η 값 레이블
\node[below, font=\footnotesize] at (0, -0.15) {$0$};
\node[below, font=\footnotesize\bfseries] at (3.2, -0.15) {$0.8$};
\node[below, font=\footnotesize\bfseries] at (6.4, -0.15) {$1.0$};
\node[below, font=\footnotesize\bfseries] at (11.2, -0.15) {$1.5$};

% 영역 레이블 (상단)
\node[font=\small\bfseries, green!50!black] at (1.6, 3.0) {안정};
\node[font=\small\bfseries, yellow!60!black] at (4.8, 3.0) {주의};
\node[font=\small\bfseries, orange!60!black] at (8.8, 3.0) {경고};
\node[font=\small\bfseries, red!60!black] at (12.8, 3.0) {위기};

% 영문 등급
\node[font=\scriptsize, green!40!black] at (1.6, 2.5) {Green};
\node[font=\scriptsize, yellow!50!black] at (4.8, 2.5) {Yellow};
\node[font=\scriptsize, orange!50!black] at (8.8, 2.5) {Orange};
\node[font=\scriptsize, red!50!black] at (12.8, 2.5) {Red};

% 핵심 조건 설명
\node[font=\scriptsize, align=center] at (1.6, 1.5)
  {관측성에 의한\\$\DRTO$ 단축\\$E_{O,bnd}$ 지배};
\node[font=\scriptsize, align=center] at (4.8, 1.5)
  {불확실성 증가\\관측성 효과 상쇄\\$E_{O,bnd} + E_{U,bnd} \approx 0$};
\node[font=\scriptsize, align=center] at (8.8, 1.5)
  {$\DRTO > R_0$\\불확실성 프리미엄\\C2 기반 자원 추가};
\node[font=\scriptsize, align=center] at (12.8, 1.5)
  {파레토 활성화\\$0.52$배 감쇠\\BCP 전면 가동};

% R₀ 기준선 표시
\draw[dashed, blue!60, thick] (6.4, -0.5) -- (6.4, -0.9)
    node[below, font=\scriptsize, blue!60] {$\eta = 1.0$: $\DRTO = R_0$ 기준};

\end{tikzpicture}
}% end resizebox
\caption{$\eta$ 기반 에스컬레이션 임계값 \m{다이어그램(Green / Yellow / Orange / Red)}}
\label{fig:ch6-escalation-zones}
\end{figure}

에스컬레이션 임계값의 설정 근거는 다음과 같다.
$\eta = 0.8$은 C1 조건에서의 평균 $\eta = 0.853$보다 낮은 값으로,
관측성의 단축 효과가 충분히 발현되어 추가 대응이 불필요한 안정 상태를 의미한다.
$\eta = 1.0$은 $\DRTO = R_0$, 즉 DRTO가 기준 RTO와 동일한 지점으로,
이 이상에서는 불확실성 프리미엄에 의해 복구 시간이 기준을 초과한다.
$\eta = 1.5$는 $\DRTO = 1.5 \times R_0 = 9.0$\,h에 대응하며,
이 이상에서는 시그모이드 바운딩의 포화 영역에 진입하여
위기관리 모드 전환이 필요하다.
이러한 단계별 에스컬레이션은 SRE의 에러 예산 소진율 기반
인시던트 대응 체계~\citep{beyer2016sre}와 유사한 구조를 가지나,
$\eta$가 관측성과 불확실성의 결합 효과를 단일 지표로 포착한다는 점에서
BCM 맥락에 최적화된 확장이다.

\chg{나아가 2단계 프레임워크와의 복합 에스컬레이션이 가능하다.}
$\eta$ 기반 4단계 등급은 \subsubsecref{subsec:two-stage-framework}의
2단계 프레임워크(Core C2 + Tail C3)와 결합하여 더욱 정교한 대응이 가능하다.
동일한 $\eta$ 등급이라도 불확실성 분포가 가우시안(C2)인지
파레토(C3)인지에 따라 위험의 성격이 질적으로 다르기 때문이다.
예를 들어, $\eta = 1.3$(Orange 등급)이 Core 영역(C2)에서 관측되면
가우시안 불확실성의 일시적 증가로 해석되어 추가 자원 투입과 모니터링 강화가
1차 대응이 된다. 그러나 동일 $\eta$가 Tail 영역(C3)에서 관측되면
파레토 꼬리 위험이 활성화된 상태이므로, $k_O$ 계수의 $0.52$배 감쇠로 인해
관측성 효과가 제한적이며, CVaR$_{99}$가 $+24.0\%$ 높은 극단 시나리오에
대비한 선제적 BCP 점검이 필요하다.

이를 운영 규칙으로 정리하면 다음과 같다.
\chg{먼저 Core 영역(P85.8 이하)에서는 $\eta$ 등급에 따른 표준 에스컬레이션 절차를 적용하고,
Tail 영역(P85.8 초과)에서는 동일 $\eta$ 등급이라도 한 단계 높은 대응 수준을
적용하며(예: Tail에서의 Orange $\to$ Red 수준 대응),
Core에서 Tail로의 조건 전환(C2 $\to$ C3) 자체를
에스컬레이션 트리거로 취급한다.}
이러한 복합 에스컬레이션은 $\eta$라는 단일 축의 한계를 보완하여,
불확실성의 분포 특성까지 반영한 입체적 위기 대응을 가능하게 한다.

\chg{또한 국가 위기경보 체계와도 대응된다.}
본 연구의 $\eta$ 기반 에스컬레이션 4단계는 ``재난 및 안전관리 기본법''에 근거한
\textbf{국가 위기경보 4단계}(관심--주의--경계--심각)와 구조적으로 대응된다.
\p{[}표~\vref{tab:eta-crisis-alert-mapping}\p{]}은 양 체계의 매핑을 보여준다.

\begin{table}[htbp]
\centering
\caption{$\eta$ 기반 에스컬레이션과 국가 위기경보 4단계 대응}
\label{tab:eta-crisis-alert-mapping}
\small
\begin{tabularx}{\textwidth}{c c c c c X}
\toprule
\textbf{$\eta$ 범위} & \textbf{DRTO 등급} & \textbf{색상} & \textbf{위기경보} & \textbf{색상} & \textbf{대응 원칙의 공통점} \\
\midrule
$\eta < 0.8$ & 안정 & Green & 관심 & 파랑 &
  징후 감지 수준; 정상 운영 유지, 모니터링 강화 \\
\addlinespace
$0.8 \leq \eta < 1.0$ & 주의 & Yellow & 주의 & 노랑 &
  위험 가능성 증가; 자원 점검, 갱신 주기 단축 \\
\addlinespace
$1.0 \leq \eta < 1.5$ & 경고 & Orange & 경계 & 주황 &
  기준 초과 확인; 추가 자원 투입, 이해관계자 통보 \\
\addlinespace
$\eta \geq 1.5$ & 위기 & Red & 심각 & 빨강 &
  총력 대응; BCP 전면 가동, 경영진 의사결정 \\
\bottomrule
\end{tabularx}
\end{table}

두 체계는 \chg{위험 강도에 따른 단계적 격상(graduated response),
색상코드 기반 직관적 커뮤니케이션,
각 단계별 보고 대상 및 대응 수준의 명확한 구분}이라는
공통 설계 원칙을 공유한다.
차이점은 국가 위기경보가 질적 판단에 기반하는 반면,
$\eta$ 기반 에스컬레이션은 $\DRTO / R_0$라는 \textbf{정량적 지표}에
의해 자동으로 등급이 결정된다는 점이다.
이러한 정량적 트리거는 주관적 판단 편향을 배제하고
일관된 의사결정을 지원하며, 기존 위기경보 체계에 보완적으로
통합될 수 있다.


  
%% --- 선순환 구조 ---
\subsubsection{\syncr{진단에서 향상, 대응, 환류로 이어지는 선순환 구조}}
\label{subsubsec:virtuous-cycle}

$\eta$의 \chg{평상시와 위기시} 활용이 만드는 핵심 가치는 \textbf{선순환}이다.
평상시의 성숙도 진단과 역량 향상이 위기시의 대응력을 높이고,
위기시의 에스컬레이션 이력\chg{(}Green 체류 비율,
복귀 시간, Red 도달률 등\chg{)}이 평상시의 성숙도 개선 데이터로 환류된다.
이 순환의 공통 언어가 $\eta$이며,
BCMS가 ``문서 기반 형식적 체계''에서 ``데이터 기반 실행적 체계''로
전환되는 것이 $\DRTO$ 프레임워크가 실무에 제공하는 핵심 가치이다.


\endgroup

\subsubsection{\syncr{실무 적용 가능성 예시}}
\label{subsubsec:practical-application-example}

\syncr{지금까지 제시한 $\eta$ 기반 \chg{평상시와 위기시} 이중 활용 체계가 실제 조직에
적용될 경우, 관측성($O_{\text{raw}}, k_O$)과 불확실성($U$)은 APM/SIEM 로그,
장애 이력 DB, CVE 공개 DB, 자산 대장 등 현장의 기존 운영 시스템에서
직접 관측되는 원시 지표로부터 산출 가능할 것이다.
참고를 위해 아래 \subsecref{app:operationalization}에 제1계층 하위 지표의
조작화 예시와 가상 기업의 4단계 라이프사이클
(진단→향상→대응→환류) 적용 가능 시나리오를 제시한다.
본 예시는 실무 적용 방향성을 보여주기 위한 참고용이며, 하위 지표의
구체적 \chg{구성과 가중치, 측정 도구}의 실증적 검증은 본 연구의 범위를 벗어나는
후속 연구 과제로 남겨둔다(\secref{sec:limitations} 참조).}



%% [본문 이관] 구 부록 J: 하위 변수 조작화 및 가상 기업 적용 예시 (참고용)
\subsection{\syncr{하위 변수 조작화 및 가상 기업 적용 예시}}
\label{app:operationalization}

\syncr{\subsubsection{하위 변수 조작화의 성격과 한계}}

\syncr{\textbf{하위 변수 조작화는 $\DRTO$ 프레임워크의 실무 적용 방향성을 보여주기 위한
참고용 예시이다.} 제시된 하위 지표의 구성, 가중치, 측정 도구 매핑,
가상 기업 수치는 \textbf{실증적으로 검증되지 않았으며}, 중견 제조업 일반
특성에서 추정한 \textbf{조건적 시나리오}(illustrative case~\citep{yin2018case})에
해당한다. 이를 규범적 가이드라인 또는 검증된 방법론으로 해석하지 말 것을 당부한다.}

\syncr{본문에서 제시한 $\DRTO$ \chg{모델과 시뮬레이션 결과, 전문가 검증}은 본 연구의
공식적 기여이나, 하위 변수 조작화와 시나리오는 \textbf{향후 연구를 위한
탐색적 제안}에 그친다. 각 하위 지표의 심리측정학적 타당도, 가중치의
산업별 튜닝, 원시 지표의 측정 도구별 \chg{정확도와 편향}, 현장 데이터 기반
\chg{수렴 및 예측} 타당도는 본 연구의 범위를 벗어나며 후속 연구에서 독립적으로
검증되어야 한다.}

\syncr{\subsubsection{조작화의 필요성과 정지점 원칙}}

\syncr{$\DRTO$ 모델의 주 변수인 관측성($O_{\text{raw}}$, $k_O$)과 불확실성($U$)은
수식적으로는 정의되어 있으나, 실제 조직에서 이들을 어떻게 산출할 것인가는
추가 조작적 정의가 필요하다. 조작화는 원칙적으로 하위 변수의 하위 변수로
무한 분해될 수 있으므로, 어느 수준에서 정지할 것인가의 \textbf{정지점
(stopping point)} 문제가 발생한다.}

\syncr{하위 변수 조작화는 측정이론의 고전적 원칙을 따라, 해당 하위 지표가 현장의
\textbf{기존 운영 시스템}(APM 도구·SIEM 도구·장애 이력 데이터베이스·CVE 공개 데이터베이스·자산 대장 등)\textbf{에서 직접 \chg{관측되고 기록되는} 원시값}에
도달하면 그 이하의 분해는 본 연구의 범위 밖으로 간주한다. 이는
\synct{bagozzi1998general}의 ``observable primitive'' 정지 규칙,
\synct{bollen1989structural}의 effect indicator 이론, 그리고
\nrev{ISO/IEC 25010:2011}의 원자 지표 표준화 원칙에 부합한다~\citep{iso25010}.}

\syncr{\chg{다만} 본 정지 규칙은 일반적 원칙을 따른 것이며,
개별 원시 지표의 심리측정학적 타당도는 측정 도구별로 별도 검증이 필요하다.}

\syncr{\subsubsection{관측성 $O_{\text{raw}}$ 조작화 예시}}

\syncr{관측성 원시값 $O_{\text{raw}} \in [0,1]$는 응용 성능 모니터링(APM)과 보안 \chg{정보 및 이벤트}
관리(SIEM)의 두 차원으로 분해 가능하다.}

\syncr{\chg{APM 차원은 $O_{\text{APM}} = w_1 R_{\text{avail}} + w_2 (1-R_{\text{error}}) + w_3 R_{\text{perf}}$로 정의한다.}}

\syncr{\chg{SIEM 차원은 $O_{\text{SIEM}} = v_1 (1-E_{\text{high}}) + v_2 (1-T_{\text{mttr}})$로 정의한다.}}

\syncr{\chg{두 차원의 합성은 $O_{\text{raw}} = \alpha \cdot O_{\text{APM}} + (1-\alpha) \cdot O_{\text{SIEM}}$이며, 본 예시에서 $\alpha = 0.7$로 가정한다.}}

\syncr{제1계층 원시 지표는 \tabref{tab:app-operationalization}에 정리하였다.}

\begin{table}[htbp]
\centering
\caption{\syncr{\chg{관측성과 불확실성, $k_O$}의 1차 조작화 \m{예시(참고용)}}}
\label{tab:app-operationalization}
\small
\syncr{\begin{tabularx}{\textwidth}{l l X l}
\toprule
\textbf{주 변수} & \textbf{1차 차원} & \textbf{원시 \m{지표(정지점)}} & \textbf{측정 출처} \\
\midrule
\multirow{5}{*}{$O_{\text{raw}}$} &
\multirow{3}{*}{$O_{\text{APM}}$} &
  $R_{\text{avail}}$: 가용시간/총시간 & APM 도구 로그 \\
& & $R_{\text{error}}$: 오류건수/총요청수 & APM 도구 로그 \\
& & $R_{\text{perf}}$: 목표응답/평균응답 & APM 도구 로그 \\
\cmidrule(l){2-4}
& \multirow{2}{*}{$O_{\text{SIEM}}$} &
  $E_{\text{high}}$: 미대응 고위험 이벤트 비율 & SIEM 도구 \\
& & $T_{\text{mttr}}$: 평균 \erf{탐지시간}{복구시간}(정규화) & SIEM 도구 \\
\midrule
\multirow{2}{*}{\jrev{$U_{\text{raw}}$}} &
  $U_{\text{hist}}$ & $\ln(P99/\jrev{\overline{T}})$ of 복구시간 & 장애 이력 DB \\
& $U_{\text{ext}}$ & \jrev{CVE CVSS 평균} & NVD/MITRE ATT\&CK \\
\midrule
$k_O$ & (단일) & 모니터링 커버리지율 & 자산 대장 \\
\bottomrule
\end{tabularx}}
\end{table}

\syncr{\chg{다만} 가중치 $w_i$, $v_i$, $\alpha$는 본 예시에서
등가중 또는 0.7로 임의 가정하였으며, \chg{산업별과 조직별} 최적값은 향후 연구에서
실증 검증되어야 한다.}

\syncr{\subsubsection{불확실성 \imod{$U$}{$U_{\text{raw}}$} 조작화 예시}}

\syncr{불확실성 \jrev{$U_{\text{raw}}$}는 과거 장애 이력 기반 $U_{\text{hist}}$와 외부 위협 지수 기반
$U_{\text{ext}}$의 두 차원으로 분해 가능하다\jrev{(두 하위 지표 모두 $[0, \infty)$
범위로 본 연구 모델의 $U_{\text{raw}}$ 정의와 일치한다)}.}

\syncr{\chg{과거 장애 기반 지표는 $U_{\text{hist}} = \ln\!\left(P99_{\text{recover}} / \overline{T}_{\text{recover}}\right) \jrev{\in [0, \infty)}$로, 복구시간 분포의 꼬리 두께를 로그 스케일로 정량화하며 장애 이력 DB에서 직접 산출한다.}}

\syncr{\chg{외부 위협 기반 지표는 $U_{\text{ext}} = \text{CVSS}_{\text{avg}} \jrev{\in [0, \infty)}$로, 최근 12개월 관련 CVE의 CVSS 점수 평균이며 NVD/MITRE ATT\&CK 공개 DB에서 산출한다.} \jrev{표준 CVSS v3.1은 $[0, 10]$이지만 조직 내부 위협 가중치 적용 시 10 초과 가능하다.}}

\syncr{\chg{두 차원의 합성은 \jrev{$U_{\text{raw}}$} $= \beta \cdot U_{\text{hist}} + (1-\beta) \cdot U_{\text{ext}}$이며, 본 예시에서 $\beta = 0.6$으로 가정한다.}}

\syncr{\chg{다만} 조직별 이력 DB 구축 수준이 상이하므로 $U_{\text{hist}}$
산출의 신뢰도는 조직마다 다를 수 있다. 신설 조직 또는 이력 부족 조직의
경우 $\beta$를 낮추거나 업종 평균을 사용하는 보정이 필요하다. \jrev{또한
$U_{\text{hist}}$와 $U_{\text{ext}}$는 각각 독립적으로 가우시안 또는 파레토
분포를 따를 수 있으며(조합 시 $2 \times 2 = 4$가지 분포 시나리오 가능),
본 연구 조건 C2/C3는 양 극단(전량 가우시안 / 전량 파레토)만 다룬다.
혼합 분포 시나리오의 실증은 향후 연구 과제이다.}}

\syncr{\subsubsection{관측성 가중치 $k_O$ 조작화 예시}}

\syncr{관측성 가중치 $k_O \in [0,1]$는 조직이 관측성에 얼마나 투자했는지를
나타내는 감도 계수로, 본 예시에서는 \textbf{모니터링 커버리지율}로 다음과 같이 단순화한다\chg{.}}

\syncr{$k_O = \frac{\text{감시 대상 자산 수}}{\text{전체 핵심 자산 수}}$}

\syncr{자산 대장에서 직접 산출되는 원시값이며, 감시 대상 자산은 APM/SIEM 도구에
등록된 엔드포인트 또는 프로세스를 기준으로 한다.}

\syncr{\chg{다만} 커버리지율은 ``감시의 폭''만을 측정하며, ``감시의
깊이''(로그 수집 상세도, 알람 정확도)는 별도 지표로 보완되어야 한다.
본 예시는 최소 구성만 제시한다.}

\syncr{\subsubsection{가상 기업 라이프사이클 적용 시나리오}}

\syncr{``(주)한빛정밀''(가상 기업명, 동명의 실제 기업과 무관)은
수도권 산업단지에 소재한 중견 제조업체(반도체 부품 정밀가공, 직원 420명,
연매출 1,200억 원 규모)로 가정하여 다음과 같이 시나리오를 작성한다. \nrev{ISO 22301:2019} 인증 준비 단계이며,
APM(Datadog)과 SIEM(Splunk)을 \chg{구축하여 운영한다}.}

\syncr{\chg{1단계인 평상시 진단(T-365일)에서는} 연간 BCMS 성숙도 평가 시점의 원시 지표를 다음과 같이
$R_{\text{avail}} = 0.994$, $R_{\text{error}} = 0.015$, $R_{\text{perf}} = 0.82$,
$E_{\text{high}} = 0.08$, $T_{\text{mttr}} = 0.65$, 커버리지율 $= 0.72$로 산출하였다.}

\syncr{합성 결과는 다음과 같다\chg{.} $O_{\text{APM}} = 0.933$, $O_{\text{SIEM}} = 0.785$,
$O_{\text{raw}} = 0.889$, $k_O = 0.72$, $U = 0.55$로 산출되어
$\eta \approx 1.05$ (Level 2 Repeatable) 판정된다. 모니터링 커버리지 확대 및
미대응 이벤트 축소가 개선 필요 영역으로 도출된다.}

\syncr{\chg{2단계인 평상시 개선 활동(T-365일부터 T-1일)에서는} 분기별 복구 훈련 도입, 모니터링 자산 확대(커버리지율 0.72 → 0.88),
SIEM 튜닝($E_{\text{high}}$ 0.08 → 0.03), 장애 이력 DB 정비를 수행하였다.
금년 재평가\chg{에서} $\eta \approx 0.82$\chg{를 달성하여} \textbf{\erf{Level 4 Managed}{Level 3 Defined}}로 \chg{향상되었다}.
월별 $\eta$ 추이를 경영진 대시보드에 정착한다.}

\syncr{\chg{3단계인 위기 발생 및 에스컬레이션(T+0h--T+6h)의} 가상 시나리오는 생산라인 제어 PLC 장애와 보조 전원계통 이상이 발생하였고
(인명 피해 없음, 설비 손상만 발생), 시간대별 $\eta$ 추이와 등급 변화는
\tabref{tab:app-lifecycle-eta}에 정리한다.}

\begin{table}[htbp]
\centering
\caption{\syncr{가상 기업 4단계 라이프사이클 $\eta$ \m{추이(참고 시나리오)}}}
\label{tab:app-lifecycle-eta}
\small
\syncr{\begin{tabularx}{\textwidth}{c l X c c}
\toprule
\textbf{시각} & \textbf{단계} & \textbf{상황} & \textbf{$\eta$} & \textbf{등급} \\
\midrule
$T{-}365$d & 1단계 진단 & 초기 \m{상태(Level 2)} & 1.05 & \erf{Yellow}{Orange} \\
$T{-}1$d & 2단계 개선 후 & 모니터링 확대 + SIEM \m{튜닝(\erf{Level 4}{Level 3})} & 0.82 & \erf{Green}{Yellow} \\
\midrule
$T{+}0$h & 3단계 정상 & 사건 직전 & 0.82 & \erf{Green}{Yellow} \\
$T{+}1$h & 3단계 탐지 & PLC 장애 탐지 & 1.15 & Orange \\
$T{+}2$h & 3단계 확산 & 전원계통 이상 확산 & 1.48 & Orange \\
$T{+}3$h & 3단계 정점 & 복구 불확실성 급등 & 1.62 & \textbf{Red} \\
$T{+}4$h & 3단계 전환 & 외부 지원 투입 & 1.45 & Orange \\
$T{+}6$h & 3단계 부분복구 & 주요 지표 안정화 & 1.08 & \erf{Yellow}{Orange} \\
\midrule
$T{+}72$h & 4단계 환류 & 사후 분석 후 재평가 & 0.84 & \erf{Green}{Yellow} \\
\bottomrule
\end{tabularx}}
\end{table}

\syncr{에스컬레이션\chg{은} $T{+}1$h(Orange)\chg{에} \chg{복구팀과 부서장} 자동 통지\chg{로 발동되어,}
$T{+}3$h(Red) 경영진 \chg{긴급 보고와 BCP 전면 가동을 거쳐,} $T{+}4$h 국가 위기경보 체계
연동 검토로 단계화되며 5.3.3.3(위기 시 정책 활용 방안)에서 상세 내용을 참조한다.}

\syncr{\chg{4단계인 복구와 환류(T+24h--T+72h)의} 사후 분석에서는 실제 복구 시간 $T_{\text{actual}} = 8.2$h vs $\DRTO$ 예측 7.6h
→ 상대 오차 $E_r = 7.3\%$ (수용 범위)이다. $O(t)$ 사각지대로 전원계통 센서
커버리지 부족을 식별하고, 성숙도 ``Managed → Optimized'' 이행 로드맵에
반영한다. 이번 사건을 $U_{\text{hist}}$ DB에 환류하여 내년도 진단의 입력으로 사용한다.}

\syncr{\chg{다만} 본 시나리오의 모든 수치는 중견 제조업 일반 특성에서
추정한 \textbf{조건적 예시}이며, 실제 기업의 실적 데이터가 아니다.
개별 조직의 적용 시 파라미터 재산출과 민감도 분석이 필수적이다.}

\syncr{\subsubsection{원시 지표별 측정 도구 매핑 예시}}

\syncr{\chg{앞의 관측성 $O_{\text{raw}}$, 불확실성 $U_{\text{raw}}$, 가중치 $k_O$ 조작화 예시에서 정의한} 원시 지표별 대표 측정 도구 매핑은
\tabref{tab:app-tool-mapping}에 정리하였다. 본 매핑은 \textbf{참고용}이며,
특정 도구의 추천이 아니다.}

\begin{table}[htbp]
\centering
\caption{\syncr{원시 지표별 측정 도구 매핑 \m{예시(참고용, 도구 권장 아님)}}}
\label{tab:app-tool-mapping}
\small
\syncr{\begin{tabularx}{\textwidth}{l l X}
\toprule
\textbf{원시 지표} & \textbf{측정 도구 유형} & \textbf{대표 \m{도구(참고)}} \\
\midrule
$R_{\text{avail}}$, $R_{\text{error}}$, $R_{\text{perf}}$ & APM & Datadog, New Relic, Prometheus+Grafana, Dynatrace \\
\addlinespace
$E_{\text{high}}$, $T_{\text{mttr}}$ & SIEM & Splunk, IBM QRadar, ArcSight, Elastic SIEM \\
\addlinespace
$U_{\text{hist}}$ & 장애 이력 DB & ServiceNow IRM, Jira Service Management, 자체 ITSM \\
\addlinespace
$U_{\text{ext}}$ & 위협 인텔리전스 & NVD, MITRE ATT\&CK, Recorded Future, CrowdStrike \\
\addlinespace
$k_O$ & 자산 대장/CMDB & ServiceNow CMDB, Device42, Lansweeper \\
\bottomrule
\end{tabularx}}
\end{table}

\nrev{ISO/IEC 25010:2011} \syncr{품질 모델의 Reliability(신뢰성),
Performance Efficiency(성능), Security(보안) 하위 속성이 각 원시 지표와
자연스럽게 대응되므로, 국제 표준 품질 지표 프레임워크와의 연계가 가능하다.}

\syncr{\chg{다만} 상기 도구는 예시이며, 조직의 기존 \chg{인프라와 예산, 규제}
환경에 따라 대체 선택이 필요하다.}

\syncr{\subsubsection{한계 종합 및 향후 실증 과제}}

\syncr{하위 변수 조작화와 시나리오는 참고용 예시이며, 다음 네 영역에서
후속 실증 연구가 필요하다.}

\syncr{\chg{첫째, 가중치 실증 튜닝이다.} $w_i$, $v_i$, $\alpha$, $\beta$를 \chg{산업별·조직 규모별}로 조정하는 다중 회귀 또는 최적화 연구가 필요하며,
본 예시는 등가중 또는 임의 고정값을 사용하였다.}

\syncr{\chg{둘째, 측정 도구 신뢰도이다.} 각 원시 지표의 측정 도구별 \chg{정확도와 편향}
검증이 필요하며, APM/SIEM 도구별 측정 알고리즘이 상이하므로
교차 타당도(cross-validity) 연구가 선결 과제이다.}

\syncr{\chg{셋째, 내용 타당도 재해석이다.} 본 연구의 \jrev{H9} 전문가 검증
(Section C 파라미터 적정성 6문항, S-CVI/Ave $= 0.985$)은 조작화의 내용
타당도의 간접 근거로 재해석 가능하나, 조작화 구성 자체에 대한 독립적
내용 타당도 검증은 별도 연구가 필요하다.}

\syncr{\chg{넷째, 수렴 및 예측 타당도이다.} 실제 조직 장애 이력과 $\eta$ 예측값의
상관 분석(종단 연구)을 통해 \chg{수렴, 판별, 예측} 타당도를 확보하는 실증 연구가
필요하다. 이는 본 연구 5.4.1 ``시뮬레이션 기반 검증의 한계''와 5.4.4
``시간 동역학의 부재''에서 제시한 향후 연구 방향과 연결된다.}

\syncr{종합적으로 하위 변수 조작화는 $\DRTO$ 프레임워크의 실무 적용 \textbf{방향성}을
제시하였으나, 본격 실무 배포 이전에는 상기 네 영역의 실증 연구가 선결되어야
한다.}



\FloatBarrier


% ===================================================================
% 5.3.4 산업별 적용 가이드라인 [구 5.3.2]
% ===================================================================
\subsection{\chg{산업별 적용 가이드라인}}
\label{subsec:practical-industry}

\begin{p5add}
\chg{본 항은 $\DRTO$ 모델의 산업별 적용 매트릭스를 제공한다. 먼저 P85.8 교차점에 기반한 2단계
프레임워크(Core C2 + Tail C3)로 일상 운영과 극단 시나리오의 분리 운용 원칙을 정립하고, 이어 6개 대표
산업(제조, 금융, 의료, IT 서비스, 물류, 소매)에 대한 $R_0$/$R_{\max}$/$\kappa$ 설정값과 적용 사례를
매트릭스 형태로 제시한다. 산업별 권장은 앞 항들의 회귀식 산출(\subsecref{subsec:practical-regression-tool})과
이중채널 자원 배분(\subsecref{subsec:practical-dualchannel-tool})을 각 산업의 분포 특성에 적용한 결과이다.}
\end{p5add}

\begingroup
  \renewcommand{\section}[1]{}
  \renewcommand{\subsection}[1]{\subsubsection{\syncr{#1}}}

  %% -------------------------------------------------------------------
\subsection{\vrev{Core(C2)와 Tail(C3)의 2단계 프레임워크}}
\label{subsec:two-stage-framework}
%% -------------------------------------------------------------------

앞서 논의한 교차점 분석과 이중채널 감쇠 결과에 기반하여,
일상 운영과 극단 시나리오를 분리하는 2단계(two-stage) $\DRTO$ 운영 프레임워크를
제안한다. \p{[}그림~\vref{fig:ch6-implementation-framework}\p{]}는 모니터링에서 에스컬레이션까지의
전체 운영 흐름을 도식화한 것이다.
Core(C2)와 Tail(C3)의 2단계 프레임워크에서는 2단계 분리의 원리와 Core/Tail 운영 전략을 중심으로 기술하며,
$\eta$ 기반 에스컬레이션 등급 체계와 구체적 대응 절차는
\subsecref{subsec:escalation-thresholds}에서 상세히 다룬다.

\begin{figure}[htbp]
\centering
\resizebox{\textwidth}{!}{%
\begin{tikzpicture}[
    stepbox/.style={draw, rounded corners=6pt, minimum width=36mm, minimum height=20mm,
                 align=center, font=\small, line width=0.7pt},
    detail/.style={draw, rounded corners=3pt, fill=white, minimum width=34mm,
                   align=left, font=\scriptsize, inner sep=4pt},
    conn/.style={-{Stealth[length=2.5mm]}, thick, blue!60},
    feedback/.style={-{Stealth[length=2mm]}, thick, black, dashed},
]

% 4단계 박스
\node[stepbox, fill=green!12]  (s1) at (0, 0)
  {\textbf{1. 모니터링}\\$O(t)$, $U(t)$ 수집\\APM/SIEM 연동};
\node[stepbox, fill=blue!12]   (s2) at (4.5, 0)
  {\textbf{2. 계산}\\$\DRTO$ 산출\\C2 또는 C3 적용};
\node[stepbox, fill=orange!12] (s3) at (9, 0)
  {\textbf{3. 의사결정}\\$\eta$ 평가\\Core vs Tail 판별};
\node[stepbox, fill=red!12]    (s4) at (13.5, 0)
  {\textbf{4. 에스컬레이션}\\$\eta$ 기반\\단계별 대응};

% 연결
\draw[conn] (s1) -- (s2);
\draw[conn] (s2) -- (s3);
\draw[conn] (s3) -- (s4);

% 세부 내용 (피드백 루프보다 먼저 배치)
\node[detail] (d1) at (0, -2.5) {
  $\bullet$ 관측성: APM 지표 $\to$ $O \in [0,1]$\\
  $\bullet$ 불확실성: 장애 이력 $\to$ $U$\\
  $\bullet$ 30분 주기 갱신\\
  $\bullet$ 이상 탐지 시 즉시 갱신
};
\node[detail] (d2) at (4.5, -2.5) {
  $\bullet$ Core ($\leq$ P85.8): C2 적용\\
  $\bullet$ Tail ($>$ P85.8): C3 적용\\
  $\bullet$ $\DRTO = R_0 + E_{O,bnd} + E_{U,bnd}$\\
  $\bullet$ $\eta = \DRTO / R_0$ 산출
};
\node[detail] (d3) at (9, -2.5) {
  $\bullet$ $\eta < 0.8$: 안정 구간\\
  $\bullet$ $0.8 \leq \eta < 1.0$: 주의\\
  $\bullet$ $1.0 \leq \eta < 1.5$: 경고\\
  $\bullet$ $\eta \geq 1.5$: 위기
};
\node[detail] (d4) at (13.5, -2.5) {
  $\bullet$ 단계별 보고 대상 결정\\
  $\bullet$ 추가 자원 배정\\
  $\bullet$ 비상 계획 가동 판단\\
  $\bullet$ 사후 분석 및 교정
};

% 세부 연결
\draw[dashed, black, line width=0.8pt] (s1.south) -- (d1.north);
\draw[dashed, black, line width=0.8pt] (s2.south) -- (d2.north);
\draw[dashed, black, line width=0.8pt] (s3.south) -- (d3.north);
\draw[dashed, black, line width=0.8pt] (s4.south) -- (d4.north);

% 피드백 루프 (상단 우회: s4 → 위 → s1)
\draw[feedback] (s4.north) -- ++(0, 1.0) -| (s1.north)
    node[pos=0.5, above, font=\scriptsize, text=black] {파라미터 교정 피드백};

% 단계 레이블 (피드백 루프 위에 배치)
\node[font=\footnotesize\bfseries, text=black] at (0, 2.8)  {실시간 입력};
\node[font=\footnotesize\bfseries, text=black]  at (4.5, 2.8) {모델 적용};
\node[font=\footnotesize\bfseries, text=black] at (9, 2.8) {상태 판별};
\node[font=\footnotesize\bfseries, text=black]   at (13.5, 2.8) {대응 실행};

\end{tikzpicture}
}% end resizebox
\caption{모니터링에서 계산, 의사결정, 에스컬레이션으로 이어지는 $\DRTO$ 구현 프레임워크}
\label{fig:ch6-implementation-framework}
\end{figure}

2단계 프레임워크의 핵심 원리는 \p{[}표~\vref{tab:two-stage-summary}\p{]}와 같다.

\begin{table}[htbp]
\centering
\caption{Core(C2)와 Tail(C3)의 2단계 프레임워크 요약}
\label{tab:two-stage-summary}
\small
\begin{tabularx}{\textwidth}{l X X}
\toprule
\textbf{구분} & \textbf{1단계 --- Core \m{영역(C2 기반)}} & \textbf{2단계 --- Tail \m{영역(C3 기반)}} \\
\midrule
\jrev{적용 \m{영역(P85.8 분할)}}     & P85.8 \m{이하(일상적 운영)}             & P85.8 \m{초과(극단 시나리오)} \\
\addlinespace
\jrev{환경 분포}  & 가우시안 ($U \sim N(3,1)$)           & 파레토 ($U \sim 6 \cdot \mathrm{Pa}(3)$) \\
\addlinespace
\jrev{환경 전체 평균 $\bar{\eta}$}  & $\bar{\eta}_{\text{C2}} \approx 1.68$ & $\bar{\eta}_{\text{C3}} \approx 1.51$ \\
\addlinespace
\jrev{C3 내 분할 영역 $k_O$ 계수}    & \jrev{Core: $-0.798$}                             & \jrev{Tail: $-0.415$ ($0.52$배 감쇠)} \\
\addlinespace
CVaR$_{99}$ \jrev{(환경 전체)}   & $3.180$                              & $3.944$ ($+24.0\%$) \\
\addlinespace
\jrev{환경 전체 $R^2$ (9변수)} & $0.998$                              & $0.858$ \\
\addlinespace
운영 목표      & 효율적 \m{복구(관측성 효과 극대화)}       & 보수적 \m{대응(꼬리 위험 대비)} \\
\addlinespace
자원 배분      & 정상 복구 절차                        & 추가 자원 투입, 비상 계획 연계 \\
\bottomrule
\end{tabularx}
\end{table}

\jrev{[표 5-14]는 Core/Tail의 두 가지 의미를 함께 다룬다. \chg{첫째는 P85.8 기준 \textit{분할 영역}(적용 영역 행)이고, 둘째는 분할 영역에 적용되는 \textit{환경 모델 C2/C3}(환경 모델, 환경 평균, 환경 $R^2$ 행)이며, 셋째는 C3 조건 \textit{내부}의 분할 회귀 결과($k_O$ 계수 행, 4.10절 dual\_channel 결과)이다.} 각 지표의 산출 분석 대상이 다르므로 행 단위로 의미를 구분하여 해석해야 한다.}

\textbf{1단계(Core, C2 기반)}는 P85.8 이하의 일상적 운영 상황에서 적용되며,
가우시안 불확실성 분포(C2 조건)에 기반한 효율적 복구를 추구한다.
\jrev{C2 조건 전체} 평균 $\bar{\eta} \approx 1.68$이며,
\jrev{C3 조건의 분할 회귀 결과 Core 영역의 $k_O$ 계수는 $-0.798$로} 관측성의 효과가 충분히 발현된다.
\textbf{2단계(Tail, C3 기반)}는 P85.8 초과의 극단 시나리오에서 활성화되며,
파레토 불확실성 분포(C3 조건)의 특성을 반영한다.
\jrev{C3 분할 회귀 결과 Tail 영역의 $k_O$ 계수가 $-0.415$로 $0.52$배 감쇠되고},
CVaR$_{99}$가 C2 대비 $+24.0\%$ 높아져, 불확실성 관리 중심의 추가 자원
투입이 정당화된다.

이러한 2단계 분리는 앞서 규명된 P85.8 분할 기준점의 이론적 근거에
기반하며, 조직이 일상 운영에서는 효율성을, 극단 상황에서는
안정성을 각각 최적화할 수 있는 구조적 기반을 제공한다.


  
%% -------------------------------------------------------------------
\subsection{산업별 권장 파라미터}
\label{subsec:industry-guidelines}
%% -------------------------------------------------------------------

$\DRTO$ 프레임워크의 파라미터는 산업별 특성에 따라 교정(calibration)이
필요하다. \chref{ch:results}의 시뮬레이션 결과, 전문가 의견, 그리고 연구모형의 산업
시나리오 설계를 종합하여, 6개 주요 산업에 대한 권장 파라미터 설정을
\p{[}표~\vref{tab:ch6-industry-params}\p{]}에 제시한다.

\begin{table}[htbp]
\centering
\caption{6개 산업별 $\DRTO$ 권장 파라미터}
\label{tab:ch6-industry-params}
\small
\begin{tabularx}{\textwidth}{l c c c c X}
\toprule
\textbf{산업} & \textbf{$R_0$ (h)} & \textbf{$R_{\max}$ (h)} & \textbf{$O_{\text{raw}}$ 범위} & \textbf{권장 조건} & \textbf{적용 근거} \\
\midrule
제조업       & 4.0 & 16.0 & 0.50--0.70 & C2
  & 생산 설비 고장은 가우시안 분포 적합, 생산 계획 연동으로 $R_0$ 엄격 관리 필요 \\
\addlinespace
금융         & 1.0 & 4.0  & 0.70--0.85 & C3
  & 규제 준수 요건(RTO $\leq$ 2h) 엄격, 시장 변동$\cdot$사이버 공격의 파레토 특성 반영 필수 \\
\addlinespace
의료         & 2.0 & 8.0  & 0.60--0.75 & C3
  & 환자 안전이 최우선, 의료기기 장애, 감염 확산 시나리오의 파레토 위험 대비 \\
\addlinespace
IT 서비스    & 0.5 & 2.0  & 0.75--0.90 & C2
  & SLA 기반 복구, 인프라 장애의 대다수가 가우시안 근사 가능, 높은 관측성 인프라 \\
\addlinespace
물류         & 6.0 & 24.0 & 0.45--0.60 & C2
  & 공급망 연속성, 계절 변동$\cdot$물류 지연은 정규분포 근사, 가장 긴 허용 RTO \\
\addlinespace
소매         & 3.0 & 12.0 & 0.50--0.65 & C2
  & 매출 손실 최소화, 일상적 IT 장애 중심, POS$\cdot$결제 시스템 관측성 활용 \\
\bottomrule
\end{tabularx}
\end{table}

\chg{[표 5-15]의 공통 파라미터는 다음과 같다. 곡률 매개변수는 $\kappa = 1.2$로 \jrev{\citet{lee2026drto}}에서 확정된
매개변수이며, 본 연구는 \frev{분포 환경 적합성} 종합 명제로 \frev{운영 적합성을 시사한다}(\secref{sec:ch4_hypothesis_summary}
참조). 관측성 가중치는 권장 범위 $k_O \in [0.4, 0.6]$이고, C2는 가우시안 불확실성 분포, C3는 파레토
불확실성 분포이며, $R_0 / R_{\max}$은 연구모형 기본 설정으로 1:4 비율을 유지한다. $O_{\text{raw}}$ 범위는 해당 산업의
전형적 모니터링 성숙도를 반영한다.}

\p{[}그림~\vref{fig:ch6-industry-summary}\p{]}는 6개 산업의 핵심 파라미터와
권장 조건을 종합적으로 시각화한 것이다.

\begin{figure}[htbp]
\centering
\resizebox{\textwidth}{!}{%
\begin{tikzpicture}[
    indbox/.style={draw, rounded corners=5pt, minimum width=30mm, minimum height=22mm,
                   align=center, font=\small, line width=0.6pt},
    modelC2/.style={draw, rounded corners=3pt, fill=blue!10, minimum width=12mm,
                    minimum height=6mm, font=\scriptsize\bfseries},
    modelC3/.style={draw, rounded corners=3pt, fill=red!10, minimum width=12mm,
                    minimum height=6mm, font=\scriptsize\bfseries},
    conn/.style={-{Stealth[length=2mm]}, thick, black},
]

% 산업 박스
\node[indbox, fill=green!8]  (mfg) at (0, 0)
  {\textbf{제조업}\\$R_0 = 4.0$h\\$O = 0.50$--$0.70$};
\node[indbox, fill=blue!8]   (fin) at (3.2, 0)
  {\textbf{금융}\\$R_0 = 1.0$h\\$O = 0.70$--$0.85$};
\node[indbox, fill=red!8]    (hc)  at (6.4, 0)
  {\textbf{의료}\\$R_0 = 2.0$h\\$O = 0.60$--$0.75$};
\node[indbox, fill=purple!8] (it)  at (9.6, 0)
  {\textbf{IT 서비스}\\$R_0 = 0.5$h\\$O = 0.75$--$0.90$};
\node[indbox, fill=orange!8] (log) at (12.8, 0)
  {\textbf{물류}\\$R_0 = 6.0$h\\$O = 0.45$--$0.60$};
\node[indbox, fill=yellow!10](ret) at (16.0, 0)
  {\textbf{소매}\\$R_0 = 3.0$h\\$O = 0.50$--$0.65$};

% 모델 레이블
\node[modelC2] at (0, -1.7) {C2};
\node[modelC3] at (3.2, -1.7) {C3};
\node[modelC3] at (6.4, -1.7) {C3};
\node[modelC2] at (9.6, -1.7) {C2};
\node[modelC2] at (12.8, -1.7) {C2};
\node[modelC2] at (16.0, -1.7) {C2};

% 연결
\draw[dashed, black] (mfg.south) -- (0, -1.4);
\draw[dashed, black] (fin.south) -- (3.2, -1.4);
\draw[dashed, black] (hc.south)  -- (6.4, -1.4);
\draw[dashed, black] (it.south)  -- (9.6, -1.4);
\draw[dashed, black] (log.south) -- (12.8, -1.4);
\draw[dashed, black] (ret.south) -- (16.0, -1.4);

% R₀ 크기 축 (상단)
\draw[thick, -{Stealth[length=2.5mm]}, black] (-1.2, 2.0) -- (17.2, 2.0)
    node[right, font=\scriptsize, black] {$R_0$ 증가};
\node[font=\scriptsize, black, above] at (3.2, 2.0) {짧은 RTO};
\node[font=\scriptsize, black, above] at (12.8, 2.0) {긴 RTO};

% 범례
\node[modelC2] at (5.5, -2.8) {C2};
\node[font=\scriptsize, anchor=west] at (6.2, -2.8) {가우시안 (일상적 장애)};
\node[modelC3] at (10.5, -2.8) {C3};
\node[font=\scriptsize, anchor=west] at (11.2, -2.8) {파레토 (꼬리 위험 필수)};

\end{tikzpicture}
}% end resizebox
\caption{6개 산업 $\DRTO$ 적용 요약 다이어그램}
\label{fig:ch6-industry-summary}
\end{figure}

산업별 조건 선택의 핵심 기준은 \textbf{파레토 위험의 존재 여부}이다.
금융과 의료는 사이버 공격, 시장 급변, 감염 확산 등 극단적 사건이
복구 시간을 급격히 연장할 수 있으므로 C3(파레토)를 권장한다.
C3의 CVaR$_{99}$가 C2 대비 $+24.0\%$ 높은 점은 이러한 산업에서
꼬리 위험을 사전에 반영해야 하는 정량적 근거가 된다.
반면, 제조업, IT 서비스, 물류, 소매는 장애의 대다수가 가우시안
분포로 근사 가능하므로 C2가 효율적이며, 필요 시 극단 시나리오에 한하여
C3를 보완적으로 적용할 수 있다.

분포 유형의 판단 오류는 복구 자원의 배정에서 두 가지 상반된 위험으로 이어진다.
가우시안 환경(C2)에 파레토 수준(C3)의 자원을 배정하는 과다예비(over-provisioning)는
불확실성을 과대 추정한 결과로, 유휴 대기 인력이나 미사용 예비 시스템, 과도한 예비
부품과 같은 자원이 상시 유지되어 비용 효율이 저하되나 서비스 수준 협약(SLA) 위반
위험은 낮다. 예를 들어 정상적 IT 장애 복구에 대해 랜섬웨어 대응 수준의 인력과 예비
시스템을 상시 유지하면 유휴 자원 비용이 불필요하게 증가한다. 반대로 파레토
환경(C3)에 가우시안 수준(C2)의 자원만 배정하는 과소예비(under-provisioning)는
불확실성을 과소 추정한 결과로, 비용은 절감되나 극단적 사건 발생 시 SLA 위반과
업무 중단 장기화 위험이 급증한다. 공급망 복합 장애 가능성이 있는 물류 기업이 단순
장비 교체 수준의 예비만 확보하면, P99 수준의 극단적 지연(약 $13.2$시간 이상)이
발생할 때 복구 역량이 부족하여 업무 연속성이 심각하게 위협받는다. 따라서 산업별
조건 선택은 단순한 분포 가정의 문제가 아니라 자원 배정의 비용과 위험을 균형 짓는
의사결정이며, 파레토 위험이 존재하는 산업에서 C3를 권장하는 근거가 여기에 있다.

각 산업에 대한 세부 적용 지침은 다음과 같다.

\chg{제조업의 경우,}
생산 설비 고장에 의한 업무 중단이 주요 위험이며,
관측성은 PLC(Programmable Logic Controller), SCADA, MES(Manufacturing
Execution System) 등의 기존 인프라를 활용한다.
$R_0 = 4.0$h는 일반적인 교대제 생산 라인의 복구 목표를 반영하며,
$O = 0.50$--$0.70$은 중간 수준의 모니터링 성숙도에 해당한다.

\chg{금융의 경우,}
금융감독 당국의 IT 복구 시간 규제(예: RTO $\leq$ 2시간)가 엄격하며,
\chg{사이버 공격, 시장 급변, 결제 시스템 장애}의 파레토 특성을 고려하여
C3를 권장한다. $O = 0.70$--$0.85$의 높은 관측성은 실시간 거래
모니터링 시스템(FDS 등)을 반영한다.

\chg{의료의 경우,}
환자 안전이 최우선이며, 의료기기 장애, 전자 의무기록(EMR) 시스템 중단,
감염병 확산 시나리오에서 복구 지연이 생명에 직결된다.
$R_0 = 2.0$h는 응급 의료 시스템의 전형적 복구 목표이며,
C3 적용으로 극단 시나리오에 대한 사전 대비를 강화한다.

\chg{IT 서비스의 경우,}
SLA(Service Level Agreement) 기반의 엄격한 복구 목표를 가지며,
$R_0 = 0.5$h는 클라우드/SaaS 서비스의 전형적 RTO이다.
$O = 0.75$--$0.90$의 높은 관측성은 APM, 로그 분석, 분산 추적(distributed
tracing) 등 성숙한 모니터링 인프라를 반영한다.

\chg{물류의 경우,}
공급망의 연속성이 핵심이며, 계절 변동, 운송 지연, 창고 시스템 장애가
주요 위험이다. $R_0 = 6.0$h는 본 연구의 기본 설정과 동일하며,
$O = 0.45$--$0.60$은 물류 산업의 상대적으로 낮은 IT 관측성 성숙도를 반영한다.

\chg{소매의 경우,}
POS(Point of Sale) 시스템, 재고 관리, 전자상거래 플랫폼의 가용성이
매출에 직접 영향을 미친다. $R_0 = 3.0$h, $O = 0.50$--$0.65$는
중간 규모 소매 환경의 전형적 특성을 반영하며, C2 기반의 효율적 복구를 권장한다.


\chg{이상의 산업별 권장 파라미터를 회귀식에 대입한 구체적 $\DRTO$ 산출 예시는 회귀식 기반 산출 도구를
다루는 \subsecref{subsec:practical-regression-tool}에 제시한다.}



\endgroup
\FloatBarrier
